\section{Representations and Models}
\label{sec:methods}

\subsection{Handcrafted baseline}

The baseline represents each $V_H$ sequence with seven features: length;
charged, hydrophobic, aromatic, and cysteine fractions; cysteine count; and
$N$-glycosylation-motif count. We evaluate logistic regression, random
forest~\cite{breiman2001random}, and XGBoost~\cite{chen2016xgboost} heads.
We select XGBoost as the baseline and retain a length-only logistic regression
as a simple diagnostic.

\subsection{Frozen ESM3 and downstream heads}

We represent each sequence with residue embeddings from the frozen, open-weights
\texttt{esm3-sm-open-v1}
model~\cite{evolutionaryscale2024esm3open}. ESM3 is a 1.4-billion-parameter
transformer model that learns from public protein databases. We pass each amino acid
sequence through the model, extract the 1536-dimensional representation for each
residue (excluding start-of-sequence BOS and end-of-sequence EOS tokens),
resulting in a matrix of size $L \times 1536$ for a sequence of length $L$. To
convert these variable-length representations into a fixed-size feature vector
for downstream modeling, we compare two pooling schemes: (1) simple mean pooling
across the sequence length, yielding a 1536-dimensional vector, and (2) a
concatenated representation of the mean, maximum, and standard deviation across
the sequence, yielding a 4608-dimensional vector. Downstream models fit their
preprocessing steps (standard scaling and Principal Component Analysis [PCA] for
dimensionality reduction) strictly inside each training fold to avoid target
leakage. All downstream classifier and regressor heads, scaling, and PCA steps
use the scikit-learn library~\cite{scikit-learn}.

Classification candidates are L1/L2 logistic regression (using the SAGA
solver~\cite{defazio2014saga} for L1), random forest~\cite{breiman2001random},
XGBoost, and a PCA-50 multilayer perceptron. A separately trained additive
attention-pooling classifier is an additional capacity test. Regression
candidates are Ridge, random forest~\cite{breiman2001random}, and XGBoost, with
a length-only Ridge floor. Appendix~\ref{app:hyperparameters} lists the
hyperparameters.

\subsection{OAS LoRA adaptation}

To align general-protein representations with the specific statistical
distribution of human antibodies, we perform domain adaptation using Low-Rank
Adaptation (LoRA)~\cite{hu2022lora}. LoRA freezes the base model parameters and
injects trainable low-rank decomposition matrices into the transformer's
attention and feed-forward layers, drastically reducing the parameter space. We
configure rank-16 adapters ($\alpha=32$, dropout $0.05$), training 18,874,368
parameters out of 1.42 billion total parameters (1.33\%). Adaptation uses a
self-supervised Masked Language Modeling (MLM)
objective~\cite{devlin-etal-2019-bert}: we randomly mask or corrupt 15\% of amino
acid tokens, and the model learns to predict the original residues from
the surrounding context. We train for one epoch using the AdamW
optimizer~\cite{LoshchilovH19} with an effective batch size of 32. Performance
evaluation uses cross-entropy loss and perplexity (the exponential of the
loss, representing the effective branching factor when predicting residues). Our
run achieves a clonotype-held-out validation loss of $0.5025$ and a perplexity
of $1.653$. We merge the learned adapter weights directly into the base model
parameters before extracting embeddings.

\subsection{Per-residue VAE and controls}
\label{sec:vae_methods}

To analyze sequence maturation and interpolation, we train a Variational
Autoencoder (VAE) to compress 1536-dimensional residue-level embeddings into a
smooth, lower-dimensional latent space. Variational autoencoders are generative
models consisting of an encoder that maps input vectors to a latent probability
distribution, and a decoder that reconstructs the inputs from samples drawn from
this distribution. Our VAE architecture uses fully connected layers with hidden
dimensions of 512 and 256, compressing the input to a 32-dimensional latent
variable $z$. The loss function balances reconstruction error (Mean Squared
Error) with a regularization term (Kullback-Leibler [KL] divergence) that
penalizes deviation from a standard normal prior distribution to ensure a smooth
latent space. We scale the KL penalty by $\beta=0.005$ using a 25-epoch linear
warm-up to prevent posterior collapse, wherein the model ignores the latent
codes. We fit the standard scaler on training residue vectors only.
We train separate frozen-embedding and LoRA-embedding VAEs with this same
architecture for 300 epochs.

We globally align endpoint sequences with the Needleman-Wunsch
algorithm~\cite{needleman1970general} through the Biopython
library~\cite{cock2009biopython} before residue-wise interpolation. For a UCA
latent $z_u$ and mature latent $z_m$, the VAE path is
\begin{equation} z(\alpha)=(1-\alpha)z_u+\alpha z_m,\qquad \alpha\in[0,1]. \end{equation}
We evaluate three embedding-space paths: the decoded VAE latent path; a straight line between VAE-reconstructed endpoints (\texttt{sl\_rec}), which controls for reconstruction; and a straight line between raw ESM3 endpoints (\texttt{sl\_raw}), which avoids the VAE bottleneck.

We perform four evaluations of the latent space geometry. The first evaluation (inferred intermediate search) searches an 11-step VRC34 UCA-to-mature path for the closest point to each of seven inferred intermediates (inferred ancestral sequences representing discrete evolutionary steps during maturation). The second evaluation (maturation midpoint test) tests four exact midpoints of ordered VRC34 triples. The third evaluation (maturation correlation analysis) correlates a VAE latent-axis coordinate with the sequence edit distance (the number of single-amino-acid substitutions, insertions, or deletions) from the UCA for four held-out lineages, serving as a proxy for evolutionary time. The fourth evaluation (latent space geometry check) compares VAE and reconstructed-line midpoints (geometric centers of the interpolation paths) for 75 sampled triples across 15 lineages to evaluate latent space smoothness. The maturation correlation analysis lacks a raw-embedding projection control and therefore cannot isolate a VAE or LoRA effect.
